\section{The mixed Gaussian weight-five conjecture is a theorem}
\label{sec:s4-proof}

Put $G=\beta(2)$ and
\[
 g_{a,b}=\operatorname{Im}\operatorname{Li}_{a,b}(i,1),
 \qquad
 S_p=\sum_{n=0}^{\infty}\frac{(-1)^nH_n}{(2n+1)^p}.
\]
The manuscript's equation \texttt{gauss:eq:S4-closed} was explicitly
classified as a numerically supported identity. We prove precisely that
identity. The contribution is a reduction of its mixed Gaussian term to
three specified single-color doubles and elementary products. It is not a
claim that the remaining doubles are independent or irreducible, nor a
claim that $S_4$ previously lacked every representation in special functions.
Indeed the 2021 discussion by Hintze and pisco already records the
representation by two colored double values and relates it to
hypergeometric evaluations \cite{hintze2021s4,mhzhao2020}.

\begin{theorem}[Certified mixed Gaussian reduction]\label{thm:s4}
With the series convention
$\operatorname{Li}_{a,b}(x,y)=\sum_{n>m\ge1}x^ny^m/(n^am^b)$,
\begin{align}
 S_4={}&\frac{4g_{4,1}-3g_{3,2}-9g_{2,3}}7
       +\frac{\pi^5}{224}
       -\frac{27}{224}G\zeta(3)-2\beta(4)\log2.
 \label{eq:s4-proved}
\end{align}
Equivalently,
\begin{align}
 \operatorname{Im}\operatorname{Li}_{4,1}(i,-1)
 ={}&-\frac{3g_{4,1}+3g_{3,2}+9g_{2,3}}7
       +\frac{\pi^5}{224}
       -\frac{27}{224}G\zeta(3)-2\beta(4)\log2.
 \label{eq:s4-mixed-proved}
\end{align}
The proof consists of convergent changes of variables, standard
shuffle/stuffle identities with the explicitly described single-divergence
regularization, and a finite rational linear combination supplied as an
independently replayable certificate.
\end{theorem}

\subsection{The arithmetic translation}
For $n=2k+1$,
\[
 H_k=\sum_{m=1}^{n-1}\frac{1+(-1)^m}{m}.
\]
Taking imaginary parts removes even $n$, while
$\operatorname{Im}(i^{2k+1})=(-1)^k$. Consequently, for $p\ge2$,
\begin{equation}\label{eq:s-harmonic-translation}
 S_p=\operatorname{Im}\left(
 \operatorname{Li}_{p,1}(i,1)+\operatorname{Li}_{p,1}(i,-1)\right).
\end{equation}
The interchange is justified by the majorant
$\sum_{n\ge2}H_{n-1}/n^p<\infty$.
Thus \eqref{eq:s4-proved} and \eqref{eq:s4-mixed-proved} are equivalent.

\subsection{A convergent octahedral change of variables}
Write $\omega_a=dt/(t-a)$ and
\[
 \mathcal I(a_1\cdots a_w)
 =\int_{0<t_w<\cdots<t_1<1}
     \omega_{a_1}(t_1)\cdots\omega_{a_w}(t_w),
\]
extended multilinearly in the letters. A convergent word has first letter
different from $1$ and last letter different from $0$.
For colors $c_j\in\mathbb Z/4\mathbb Z$, set
\[
 Z((s_1,c_1),\ldots,(s_d,c_d))
 =\operatorname{Li}_{s_1,\ldots,s_d}(i^{c_1},\ldots,i^{c_d}).
\]
Its integral word is
\begin{equation}\label{eq:s4-word-convention}
 Z(\boldsymbol s;\boldsymbol c)
 =(-1)^d\mathcal I\left(
 0^{s_1-1}i^{-c_1}\,
 0^{s_2-1}i^{-(c_1+c_2)}\cdots
 0^{s_d-1}i^{-(c_1+\cdots+c_d)}\right).
\end{equation}
This fixes both the color convention and every depth-dependent sign.

Consider the involution
\[
 \phi(t)=\frac{1-t}{1+t}.
\]
It preserves the puncture set $\{0,\infty,1,-1,i,-i\}$ and reverses the
real interval. Direct differentiation gives
\[
 \phi^*\omega_a=\omega_{\phi(a)}-\omega_{-1}\quad(a\ne-1),
 \qquad \phi^*\omega_{-1}=-\omega_{-1},
\]
where $\phi(0)=1$, $\phi(1)=0$, $\phi(i)=-i$, and
$\phi(-i)=i$. Denote this linear map on letters by $\tau$.
Changing variables and reversing the order of the simplex proves
\begin{equation}\label{eq:s4-convergent-duality}
 \mathcal I(a_1\cdots a_w)
 =(-1)^w\mathcal I\bigl(\tau(a_w)\cdots\tau(a_1)\bigr).
\end{equation}
No tangential-basepoint correction occurs here: an original last letter
other than $0$ transforms into first letters other than $1$, and an original
first letter other than $1$ transforms into last letters other than $0$.
Every term of the expansion therefore remains convergent.
The broader use of octahedral symmetries to supplement standard level-four
relations goes back to Zhao \cite{jqzhao2008,jqzhao2010}; the elementary
involution above is sufficient for the present certificate.

Only one instance of \eqref{eq:s4-convergent-duality} is needed. The word of
$\operatorname{Li}_{1,4}(i,1)$ is $(-i)000(-i)$, and hence
\begin{equation}\label{eq:s4-duality-seed}
 D:=\operatorname{Im}\left\{
 \mathcal I(e_{-i}e_0^3e_{-i})+
 \mathcal I\bigl((e_i-e_{-1})(e_1-e_{-1})^3(e_i-e_{-1})\bigr)\right\}=0.
\end{equation}
Here $e_a$ means the letter at $a$, multiplication is concatenation, and
powers are powers in the noncommutative word algebra. The transformed product expands into
exactly $32$ integral words, with coefficients $\pm1$. The first integral is
the single-color double; the second group is converted to colored
polylogarithms using \eqref{eq:s4-word-convention}.

\subsection{The standard relation rows}
The certificate uses two schemas, interpreted over $\mathbb Q$ before
imaginary parts are taken.
First, for two admissible multiple-polylogarithm index/color lists $u,v$,
let $\mathsf H(u,v)$ be the integral shuffle product and
$\mathsf S(u,v)$ the series stuffle product. Then
\begin{equation}\label{eq:s4-ds-row}
 R_{u,v}:=\operatorname{Im}\bigl(\mathsf H(u,v)-\mathsf S(u,v)\bigr)=0.
\end{equation}
The same row is used when $u=((1,0))$ and $v$ is admissible, with the standard
single-divergence regularization. This point must not be suppressed:
$\operatorname{Li}_1(1)$ itself is divergent. Its shuffle and stuffle
regularizations are both denoted by a formal variable $T$. A single leading
$(1,0)$ gives a polynomial of degree at most one in $T$, and the
regularization comparison is the identity in degrees zero and one.
Explicitly, its comparison operator satisfies
\[
 \rho(e^{Tu})=\exp\!\left(\sum_{n\ge2}\frac{(-1)^n\zeta(n)}n u^n\right)e^{Tu},
\]
so $\rho(1)=1$ and $\rho(T)=T$ \cite[Theorem~2.9]{jqzhao2010}.
Subtracting the two product expansions cancels their common divergent word
and leaves the convergent identity \eqref{eq:s4-ds-row}. This is the
single-divergence case of regularized double shuffle
\cite{racinet2002,jqzhao2010}. The verifier rejects all other divergent input
patterns and checks that every surviving word is admissible.

Second, for a list $z=((s_1,c_1),\ldots,(s_d,c_d))$ with
$c_j\in\{0,1\}$ and with all terms below convergent, ordinary distribution
is
\begin{equation}\label{eq:s4-distribution-row}
 Z(\boldsymbol s;2\boldsymbol c)
 -2^{w-d}\sum_{\boldsymbol\epsilon\in\{0,2\}^d}
 Z(\boldsymbol s;\boldsymbol c+\boldsymbol\epsilon)=0,
 \qquad w=s_1+\cdots+s_d.
\end{equation}
Indeed summing signs forces every nested summation index to be even, so the
factor is exactly $2^{w-d}$. The certificate multiplies some such relations
by an admissible value of complementary weight and expands that product by
integral shuffle. These are \emph{lifted convergent distribution} rows; no
regularized distribution formula is assumed. At weight one this includes
the elementary seed $\operatorname{Li}_1(-1)=
\operatorname{Li}_1(i)+\operatorname{Li}_1(-i)$.

\subsection{The exact certificate identity}
In the formal real vector space whose generators are imaginary parts of
convergent colored values of weight five, define
\begin{align}
 F={}&1120\operatorname{Im}\operatorname{Li}_{4,1}(i,-1)
       +480g_{4,1}+480g_{3,2}+1440g_{2,3}
       -1536\operatorname{Im}\operatorname{Li}_5(i)\nonumber\\
 &\quad+135\operatorname{Im}\bigl(\operatorname{Li}_2(i)
                                      \operatorname{Li}_3(1)\bigr)
       -2240\operatorname{Im}\bigl(\operatorname{Li}_4(i)
                                      \operatorname{Li}_1(-1)\bigr).
 \label{eq:s4-certificate-target}
\end{align}
The two products in this expression are expanded by integral shuffle.
Complex conjugation is imposed only by
$\operatorname{Im}Z(\boldsymbol s;-\boldsymbol c)
=-\operatorname{Im}Z(\boldsymbol s;\boldsymbol c)$.

\begin{lemma}[Finite rational certificate]\label{lem:s4-certificate}
The accompanying file \texttt{S4\_certificate.json} gives $911$ rational
numbers $q_r$ and $911$ standard-relation descriptors $R_r$ such that the
following is an exact equality of formal vectors:
\begin{equation}\label{eq:s4-certificate-equality}
 F=-1120D+\sum_{r=1}^{911}q_rR_r.
\end{equation}
The rows consist of $713$ convergent double-shuffle rows, $181$ rows with
one regularized weight-one factor, and $17$ lifted convergent distribution
rows.
\end{lemma}
\begin{proof}
The supplied program \texttt{verify\_s4.py} rebuilds every row from the
indices and colors in its descriptor. Integral shuffles are generated by
choosing interleaving positions, series stuffles by the three defining
recursive cases, and distribution by sign enumeration. The program also
reconstructs the single duality seed and the target independently of their
stored expansions. All coefficients are integers or instances of Python's
exact \texttt{Fraction} type. Summing the left side minus the right side of
\eqref{eq:s4-certificate-equality} yields the empty coefficient dictionary,
so every formal-word coefficient is exactly zero. The verifier neither
reruns a rank calculation nor uses numerical polylogarithms, integer-relation
searches, cached matrices, or modular arithmetic. Thus the certificate is
a finite rational proof of the stated formal equality.
\end{proof}

\begin{proof}[Proof of Theorem~\ref{thm:s4}]
Every $R_r$ vanishes by \eqref{eq:s4-ds-row} or
\eqref{eq:s4-distribution-row}, and $D=0$ by
\eqref{eq:s4-duality-seed}. The certificate therefore gives $F=0$.
Now
\[
 \operatorname{Im}\operatorname{Li}_5(i)=\beta(5)=\frac{5\pi^5}{1536},
 \qquad \operatorname{Li}_3(1)=\zeta(3),
 \qquad \operatorname{Li}_1(-1)=-\log2.
\]
Substituting these evaluations into \eqref{eq:s4-certificate-target} gives
\[
 1120\operatorname{Im}\operatorname{Li}_{4,1}(i,-1)
 +480g_{4,1}+480g_{3,2}+1440g_{2,3}
 -5\pi^5+135G\zeta(3)+2240\beta(4)\log2=0.
\]
Division by $1120$ proves \eqref{eq:s4-mixed-proved}, and
\eqref{eq:s-harmonic-translation} proves \eqref{eq:s4-proved}.
\end{proof}

\paragraph{Independent numerical audit.}
An independent real-quadrature evaluation uses
\[
 S_4=\frac16\int_0^1\frac{\log^3x\,\log(1+x^2)}{1+x^2}\,dx,
 \qquad
 g_{a,b}=\frac1{(a-1)!}\int_0^1(-\log x)^{a-1}
 \operatorname{Im}\frac{i\operatorname{Li}_b(ix)}{1-ix}\,dx.
\]
At $70$ decimal working digits it gives
\[
 S_4=-0.0104934562973350434533567510062017519635106188773197615131294,
\]
with discrepancy about $4.6\times10^{-71}$ between the two sides.
This numerical agreement is an audit of conventions, not a step in the
proof; the proof is the exact certificate above.

\paragraph{Scope and further work.}
The existing two relations among $g_{4,1},g_{3,2},g_{2,3},g_{1,4}$ may be
used to shorten \eqref{eq:s4-proved}; they are prior results and are not
counted as new here. The present proof establishes an upper bound on a
specific constant's expression complexity. It says nothing about a lower
bound, linear independence, or the impossibility of a classical-polylogarithm
reduction. A useful next problem is to compress the $911$ standard rows
into a short conceptual derivation, and to apply the same convergent
involution to the analogous odd-weight mixed sums $S_6,S_8,\ldots$.
